% Figure: polarisation states seen looking into the oncoming wave.
% Three panels: linear (a line), circular (a circle with rotation arrow),
% elliptical (the general ellipse).
\begin{figure}[htbp]
\centering
\begin{tikzpicture}[scale=1.0, >=latex]
% ---- Linear ----
\begin{scope}[shift={(0,0)}]
  \draw[->, enggray] (-1.4,0) -- (1.4,0) node[right, font=\scriptsize] {$E_x$};
  \draw[->, enggray] (0,-1.4) -- (0,1.4) node[above, font=\scriptsize] {$E_y$};
  \draw[engblue, very thick] (-1.1,-1.1) -- (1.1,1.1);
  \node[font=\small] at (0,-1.9) {linear};
\end{scope}
% ---- Circular ----
\begin{scope}[shift={(4,0)}]
  \draw[->, enggray] (-1.4,0) -- (1.4,0) node[right, font=\scriptsize] {$E_x$};
  \draw[->, enggray] (0,-1.4) -- (0,1.4) node[above, font=\scriptsize] {$E_y$};
  \draw[engred, very thick] (0,0) circle (1.1);
  \draw[engred, ->, very thick] (1.1,0) arc (0:60:1.1);
  \node[font=\small] at (0,-1.9) {circular};
\end{scope}
% ---- Elliptical ----
\begin{scope}[shift={(8,0)}]
  \draw[->, enggray] (-1.4,0) -- (1.4,0) node[right, font=\scriptsize] {$E_x$};
  \draw[->, enggray] (0,-1.4) -- (0,1.4) node[above, font=\scriptsize] {$E_y$};
  % ellipse rotated 30 degrees, semi-axes 1.2 and 0.6
  \draw[englime!60!black, very thick, rotate=30] (0,0) ellipse (1.2 and 0.6);
  \draw[englime!60!black, ->, very thick, rotate=30] (1.2,0) arc (0:55:1.2 and 0.6);
  \node[font=\small] at (0,-1.9) {elliptical};
\end{scope}
\end{tikzpicture}
\caption[Polarisation states of a plane wave]{Polarisation states drawn
in the plane transverse to $\vect{k}$, viewed looking into the oncoming
wave. The tip of the electric-field vector traces a line (linear
polarisation), a circle (circular polarisation), or an ellipse (the
general elliptical case). The arrow marks the sense of rotation, which
fixes the handedness. Linear polarisation is two in-phase components;
circular polarisation is two equal-amplitude components a quarter cycle
apart; the ellipse covers every intermediate amplitude and phase pair.}
\label{fig:vol04ch10:polarization-ellipse}
\end{figure}
